\section{Complete Equational Reasoning} \label{sec:pbe}

Next, we formalize our notion of \emph{equational proofs}
and present a \emph{complete} procedure for finding such
proofs.

\subsection{Preliminaries} \label{sec:pbe:prelim}
%
Figure~\ref{fig:pbe:syntax} describes the syntax of
the main ingredients.
%
\RJ{describe $\Renv$}


\begin{figure}[t!]
\begin{tabular}{>{$}r<{$} >{$}r<{$} >{$}r<{$} >{$}l<{$} >{$}l<{$} }
\emphbf{Values}       & \val   & ::=    & \vconst                                  & \mbox{constants}\\
                      &        & \spmid & \capp{\ctor}{\val}                       & \mbox{ctor-application}\\[0.05in]
\emphbf{Operators}    & \op    & ::=    & \{ =,\ \wedge,\ \ldots \}  \\[0.05in]
% \emphbf{Operators}   & \op    & ::=    & =      \spmid \leq \spmid <              & \mbox{comparisons} \\
%                      &        & \spmid & \wedge \spmid \vee \spmid \neg           & \mbox{boolean} \\[0.05in]
\emphbf{Terms} & \expr, \pred, \body
                      & ::=    & \val                                              & \mbox{values}  \\
                      &        & \spmid & x                                        & \mbox{variables} \\
                      &        & \spmid & \expr \op \expr                          & \mbox{prim-op} \\
                      &        & \spmid & \fapp{\rfun}{\expr}                      & \mbox{func-application} \\
                      &        & \spmid & \capp{\ctor}{\expr}                      & \mbox{ctor-application} \\[0.05in]
\emphbf{Functions}    & \Rfun  & ::=    & \rdef{x}{\pred}{\body}                   & \\ [0.05in]
\emphbf{Definitions}  & \Renv  & ::=    & \emptyset \spmid \rfun \mapsto \Rfun, \Renv & \\ [0.05in]
\emphbf{Environments} & \Penv  & ::=    & \ttrue \spmid \pred, \Penv               & \\ [0.05in]
\end{tabular}
\caption{{Syntax of Predicates, Terms and Reflected Functions}}
\label{fig:pbe:syntax}
\end{figure}

\mypara{Normal}
%
We say $\Penv$ is \emph{normal}
if all applications in $\Penv$ are
of the form $\fapp{f}{x}$.
%
We can easily \emph{normalize}
a $\Penv$ by introducing fresh
variables for each sub-term, and
hence in the sequel we assume
that $\pbe{\Renv}{\Penv}{\pred}$
is only invoked with normal $\Penv$.

\mypara{Stores and Inhabitation}
%
A \emph{store} $\store$ is a mapping from
(free) variables to values.
%
We write $\apply{\store}{\expr}$ for the term
obtained by substituting free variables $x$
in $\expr$ with the corresponding values in
the store $\store$.
%
We write $\satisfies{\Renv}{\store}{\expr}$
if $\trans{\Renv}{\apply{\store}{\expr}}{\ttrue}$.
%
We say $\Penv$ is \emph{inhabited}
if for some $\store$ we have
$\satisfies{\Renv}{\store}{\pred}$
for each $\pred \in \Penv$.


\mypara{Divergence}
%
We say that a closed term $\expr$
\emph{diverges under} $\Renv$ if
there is no $\val$ such that
$\trans{\Renv}{\expr}{\val}$ .

\begin{lemma}[\textbf{Divergence}] \label{lem:div} % TRIVIAL
If $\ \forall i \in \nats$ we have
$\trans{\Renv}{\expr_i}{\ctxapp{\ctx}{\expr_{i+1}}}$
then $\expr_0$ diverges under $\Renv$.
\end{lemma}

\mypara{Subterm Evaluation}
%
We write $\issubterm{\expr}{\expr'}$
if $\expr$ is a \emph{syntactic subterm}
of $\expr'$.
%
The following lemma states that
if a store satisfies a reflected
function's guard, then the evaluation
of the corresponding function body
requires evaluating \emph{every subterm}
inside the body.

\begin{lemma}\label{lem:subterm-eval}
Let $\rlookup{\Renv}{\rfun} \equiv \rdef{x}{\pred}{\body}$.
%
If $\satisfies{\Renv}{\store}{\pred_i}$ and $\issubterm{\fapp{g}{e}}{\body_i}$
then
$\trans{\Renv}{\apply{\store}{\body_i}}
              {\ctxapp{\ctx}{\rawapp{g}{\apply{\store}{\params{\expr}}}}}$.
\end{lemma}

\mypara{Totality Assumption}
%
The following assumption states that
if a store satisfies a reflected function's
guard, then the corresponding body reduces
to exactly one value.
%
Both of these are checked via
refinement typing \cite{Vazou14}.
%
\RJ{format below as "assumption" to avoid confusion}

\begin{lemma}[\textbf{Total}]\label{asm:total}
Let $\rlookup{\Renv}{\rfun} \equiv \rdef{x}{\pred}{\body}$.
%
$\rfun$ is \emph{total} if forall $\store$
%
\begin{enumerate}
\item If $\satisfies{\Renv}{\store}{\pred_i}$ then $\trans{\Renv}{\apply{\store}{\body_i}}{\val}$.
\item If $\satisfies{\Renv}{\store}{\pred_i}$, $\satisfies{\Renv}{\store}{\pred_j}$ then $i=j$.
\end{enumerate}
%
$\Renv$ is \emph{total} if every $\rfun \in \Renv$ is total.
%
\end{lemma}



\subsection{Equational Proofs}
%
Next, we formalize the notion of an equational proof.

\mypara{Instantiation}
%
A term $\predb$ is a \emph{$\Renv,\Penv$-instance} if
there exists a sub-term $\issubterm{\fapp{\rfun}{\expr}}{\Penv}$
such that:
%
\begin{itemize}
\item $\rlookup{\Renv}{\rfun} \equiv \rdef{x}{\pred_i}{\body_i}$,
\item $\smtIsValid{\Penv}{\SUBSTS{\pred_i}{x}{\expr}}$,
\item $\predb \equiv \SUBSTS{(\fapp{\rfun}{x} = \body_i)}{x}{\expr}$.
\end{itemize}
%
A set of terms $Q$ is a \emph{$\Renv, \Penv$-instance}
if every $\predb \in Q$ is an $\Renv$-instance.
%
The following properties relate instances to the exact semantics
and SMT validity.

\begin{lemma}[\textbf{Sat-Inst}] \label{lem:sat-inst}
  If   $\satisfies{\Renv}{\store}{\Penv}$ and
       $Q$ is a $\Renv,\Penv$-instance,
  then $\satisfies{\Renv}{\store}{\Penv \wedge Q}$.
\end{lemma}

\begin{lemma}[\textbf{SMT-Approx}] \label{lem:smt-approx}
If $\smtIsValid{\Penv}{\pred}$ then %forall $\store$,
   $\satisfies{\Renv}{\store}{\Penv}$ implies
   $\satisfies{\Renv}{\store}{\pred}$.
\end{lemma}

\begin{lemma}[\textbf{SMT-Inst}] \label{lem:smt-inst}
If   (1)~$\satisfies{\Renv}{\store}{\Penv}$,
     (2)~$Q$ is an $\Renv,\Penv$-instance,
     (3)~$\smtIsValid{\Penv \wedge Q}{\pred}$
then $\satisfies{\Renv}{\store}{\pred}$.
\end{lemma}



\mypara{Unfolding}
%
The \emph{unfolding} of $\Renv, \Penv$ is the
(finite) set of all $\Renv, \Penv$ instances.
%
Procedure $\uinst{\Renv}{\Penv}$ shown in
Figure~\ref{fig:uinst} computes and returns
the conjunction of $\Penv$ and the unfolding
of $\Renv, \Penv$.

\mypara{Equational Proof}
%
There is an \emph{equational proof}
that $\expr$ equals $\expr'$ under
$\Renv, \Penv$
($\eqpf{\Renv}{\Penv}{\expr}{\expr'}$)
if
$$\begin{array}{lll}
\inferrule[\rulename{Eq-Refl}]
  {\ }
  {\eqpf{\Renv}{\Penv}{\expr}{\expr}}
&
&
\inferrule[\rulename{Eq-Trans}]
  { \eqpf{\Renv}{\Penv}{\expr}{\expr''}
    \quad
    \smtIsValid{\uinst{\Renv}{\Penv \wedge \vv = \expr''}}{\vv = \expr'}
  }
  {\eqpf{\Renv}{\Penv}{\expr}{\expr'}}
\end{array}$$

\subsection{\pbesym Algorithm}


\begin{figure}[t!]
$$\begin{array}{lcl}
\toprule
\uinstsym & : & (\Renv, \Penv) \rightarrow \setof{\Pred} \\
\midrule
\uinst{\Renv}{\Penv} & = &
  \Penv
  \bigcup_{\issubterm{\fapp{\rfun}{\expr}}{\Penv}}
    \inst{\Renv}{\Penv}{\rfun}{\exprs} \\[0.1in]

%\midrule
%\instsym & : & (\Renv, \Pred, \Rfun, \params{\Expr}) \rightarrow \setof{\Pred} \\
%\midrule
\inst{\Renv}{\Penv}{\rfun}{\exprs} & = &
  \{ \SUBSTS{(\fapp{\rfun}{x} = \body_i)}{x}{\expr}
  \spmid \pred_i \Rightarrow \body_i \in \overline{d}
  ,      \smtIsValid{\Penv}{\SUBSTS{\pred_i}{x}{\expr}}
  \} \\
\quad \mbox{where} & & \\
\quad \quad \rbody{x}{d} & = & \Renv(\rfun) \\[0.1in]

\midrule
\pbesym & : & (\Renv, \Pred, \expr, \expr') \rightarrow \tbool \\
\midrule
\pbe{\Renv}{\Penv}{\expr = \expr'} & = & \pbeloop{0}{\{\Penv, \vv = \expr\}} \\
\quad \mbox{where} & & \\
\quad \quad \pbeloop{i}{\Penv_i} & & \\
\quad \quad \spmid \smtIsValid{\Penv_i}{\vv = \expr'} & = & \ttrue \\
\quad \quad \spmid \Penv' \subseteq \Penv           & = & \tfalse \\
\quad \quad \spmid \mbox{otherwise}                 & = & \pbeloop{i+1}{\Penv'} \\
\quad \quad \quad  \mbox{where}                     &   & \\
\quad \quad \quad  \quad \Penv'                     & = & \Penv \cup \uinst{\Renv}{\Penv_i} \\
\bottomrule
\end{array}$$
\caption{Proof by Evaluation}
\label{fig:pbe}
\label{fig:inst}
\label{fig:uinst}
\end{figure}

Figure~\ref{fig:pbe} shows our proof search algorithm
$\pbe{\Renv}{\Penv}{\pred}$ which takes as input
%
a set of \emph{reflected definitions} $\Renv$,
%
an \emph{hypothesis} $\Renv$, and
%
a \emph{goal} $\pred$.
%
The procedure recursively unfolds function application terms
by invoking $\uinstsym$, until either the goal can be proved
using the unfolded instances (in which case the search
returns $\ttrue$) \emph{or} no new instances are generated by
the unfolding (in which case the search returns $\tfalse$).
%
Using the Lemmas~\ref{lem:sat-inst},~\ref{lem:smt-approx},~\ref{lem:smt-inst}
that relate instantiation, SMT validity, and the exact semantics,
we can prove that the proof search is \emph{sound}:

\begin{theorem}[\textbf{Soundness}] \label{thm:pbe:sound}
If   $\pbe{\Renv}{\Penv}{\pred} = \ttrue$
then forall $\store$, we have $\satisfies{\Renv}{\store}{\Penv \Rightarrow \pred}$.
\end{theorem}

Ensuring soundness is easy: it follows from the fact
that the SMT solver reasons about a finite set of
instances in a conservative manner \ie by treating
all function applications as \emph{uninterpreted}~\cite{Nelson81}.
%

\subsection{\pbesym is Complete} \label{sec:pbe:semantics}

Next, we show that our proof search
is \emph{complete} with respect to
equational reasoning.
%
That is, if there exists an equational
proof
\eqpf{\Renv}{\Penv}{\expr}{\expr'},
then
\pbe{\Renv}{\Penv}{\expr = \expr'}
is guaranteed to returns \ttrue.
%
To prove this theorem, we introduce
the notion of \emph{bounded unfolding}
which corresponds to unfolding
definitions $n$ times.
%
We will show that unfolding preserves
congruences, and hence, that an equational
proof exists iff the goal can be proved with
\emph{some} bounded unfolding.
%
Thus, completeness follows by showing that
the proof search procedure computes the limit
(\ie fixpoint) of the bounded unfolding.
%
Finally, in \S~\ref{sec:pbe:termination}
we will show that the fixpoint is computable,
\ie there is an unfolding depth at which
we reach a fixpoint for the bounded unfolding,
and hence, that \pbesym terminates.

\mypara{Bounded Unfolding}
%
For every $\Renv, \Penv$ and $0 \leq n$,
the \emph{bounded unfolding of depth} $n$ is
defined by:
%
$$\begin{array}{lcl}
  \binst{\Renv}{\Penv}{0}     & \doteq & \Penv \\
  \binst{\Renv}{\Penv}{n+1}   & \doteq & \Penv_n \cup \uinst{\Renv}{\Penv_n} \\
  \quad \mbox{where}\ \Penv_n & =      & \binst{\Renv}{\Penv}{n}
\end{array}$$
%
That is, the unfolding at depth $n$ essentially,
performs $\uinstsym$ upto $n$ times.
%
The bounded-unfoldings yield a monotonically
non-decreasing sequence of formulas, and that
if two consecutive bounded unfoldings coincide,
then all subsequent unfoldings are the same.

\begin{lemma}[\textbf{Monotonicity}] \label{lem:bounded:mono}
$\forall 0 \leq n.\ \binst{\Renv}{\Penv}{n} \subseteq \binst{\Renv}{\Penv}{n + 1}$.
\end{lemma}

\begin{lemma}[\textbf{Fixpoint}] \label{lem:bounded:fix}
If $\Penv^* = \binst{\Renv}{\Penv}{n} = \binst{\Renv}{\Penv}{n + 1}$
then $\forall n \leq m.\ \Penv^* = \binst{\Renv}{\Penv}{m}$.
\end{lemma}

\mypara{Uncovering}
%
Next we prove that every function application
term that is \emph{uncovered} by unfolding to
depth $n$ is congruent to a term in the $n$-depth
unfolding.

\begin{lemma}[\textbf{Uncovering}] \label{lem:uncovering}
Let $\Penv_n \equiv \binst{\Renv}{\Penv, \vv = \expr}{n}$.
If   $\smtIsValid{\Penv_n}{\vv = \expr'}$
then for every $\issubterm{\fapp{\rfun}{t'}}{\expr'}$
       there exists $\issubterm{\fapp{\rfun}{t}}{\Penv_n}$
         such that $\smtIsValid{\Penv_n}{t = t'}$.
\end{lemma}

We prove the above lemma by induction on $n$
where the inductive step uses the following
property of congruence closure:

\begin{proof} (Of Lemma~\ref{lem:uncovering})
We prove the result by induction on $n$.

\smallskip{\emph{Case $n=0$:}}
Immediate as $\expr \equiv \expr'$.

\smallskip{\emph{Case $n=k+1$:}}
  \begin{align*}
  \intertext{Consider any $\expr'$ such that}
  & \smtIsValid{\Penv_{k+1}}{\vv = \expr'} \\
  \intertext{By definition $\Penv_{k+1} = \uinst{\Renv}{\Penv_k}$, hence}
  & \smtIsValid{\uinst{\Renv}{\Penv_k}, v = \expr}{\vv = \expr'} \\
  \intertext{Consider any $\issubterm{\fapp{\rfun}{t'}}{\expr'}$;
  apply Lemma~\ref{lem:congruence} to complete the proof.}
  \end{align*}
\end{proof}

\begin{lemma}[\textbf{Congruence}] \label{lem:congruence}
If   $\smtIsValid{\Penv, \vv = \expr}{\vv = \expr'}$
     and
     $\vv \not \in G, \expr, \expr'$
then for every $\issubterm{\fapp{\rfun}{t'}}{\expr'}$
       there exists $\issubterm{\fapp{\rfun}{t}}{\Penv, \expr}$
         such that $\smtIsValid{\Penv}{t = t'}$.
\end{lemma}

\begin{proof} (Of Lemma~\ref{lem:congruence})
%
Proof by induction on the structure of $\expr'$:
%
Case $\expr' \equiv x$:
  There are no subterms, hence immediate.
%
Case $\expr' \equiv c$:
  There are no subterms, hence immediate.
%
Case $\expr' \equiv \fapp{\rfun}{t'}$:
  Consider the \emph{last link} in $\Penv$
  connecting the equivalence class
  of $\vv$ (and $\expr$) to $\expr'$.
  %
  Suppose the last link is a \emph{congruence link}
  of the form $\expr = \expr'$ where
  $\expr \equiv \fapp{\rfun}{t}$ and
  $\smtIsValid{\Penv}{t = t'}$.
  %
  Then $\issubterm{\fapp{\rfun}{t}}{\Penv, \expr}$
  and we are done.
  %
  Suppose instead, the last link is an \emph{equality link}
  in $\Penv$ of the form $z = \fapp{\rfun}{t'}$.
  In this case, $\issubterm{\fapp{\rfun}{t'}}{\Penv}$
  and again, we are done.
\end{proof}

\mypara{Unfolding Preserves Equational Links}
%
Next, we use the uncovering Lemma~\ref{lem:uncovering}
to show that \emph{every instantiation} that is valid
after $n$ steps is subsumed by the $n+1$ depth unfolding.
%
That is, we show that every possible \emph{link} in
a possible equational chain can be proved equal to
the source expression via bounded unfolding.

\begin{lemma}[\textbf{Link}] \label{lem:link}
If   $\smtIsValid{\binst{\Renv}{\Penv, \vv = \expr}{n}}{\vv = \expr'}$
then $\smtIsValid{\binst{\Renv}{\Penv, \vv = \expr}{n+1}}{\uinst{\Renv}{\Penv \wedge \vv = \expr'}}$
\end{lemma}

\begin{proof} (Of Lemma~\ref{lem:link})
%
\begin{align}
\intertext{Let $G_k \doteq \binst{\Renv}{\Penv}{k}$. Let us assume that}
  & \smtIsValid{\binst{\Renv}{\Penv, \vv = \expr}{n}}{\vv = \expr'}
  \label{eq:link:1} \\
\intertext{Consider any instance}
  & \predb \equiv \fapp{\rfun}{t'} = \SUBSTS{\body_i}{x}{t'} \ \mbox{in}\ \uinst{\Renv}{\Penv \wedge \vv = \expr'}
  \label{eq:link:2} \\
\intertext{By the definition of \uinstsym, we have}
  & \issubterm{\fapp{\rfun}{t'}}{\Penv \wedge \vv = \expr'}
    \ \mbox{such that}\ \smtIsValid{\Penv}{\SUBSTS{\pred_i}{x}{t'}}
    \label{eq:link:3} \\
\intertext{By (\ref{eq:link:1}) and Lemma~\ref{lem:uncovering}
  there exists \issubterm{\fapp{\rfun}{t}}{\Penv_n} such that}
  & \smtIsValid{\Penv_n}{t = t'} \notag \\
\intertext{As $\Penv \subseteq \Penv_n$ and (\ref{eq:link:3}), by congruence}
  & \smtIsValid{\Penv_n}{\SUBSTS{\pred_i}{x}{t}} \notag \\
\intertext{Hence, the instance}
 & \fapp{\rfun}{t} = \SUBSTS{\body_i}{x}{t} \ \mbox{is in} \Penv_{n+1} \notag \\
\intertext{That is}
 & \smtIsValid{\Penv_{n+1}}{t = t' \wedge \fapp{\rfun}{t} = \SUBST{\body_i}{x}{t}} \notag \\
\intertext{And so by congruence closure}
 & \smtIsValid{\Penv_{n+1}}{\predb} \notag \\
\intertext{As the above holds for every instance, we have}
 & \smtIsValid{\Penv_{n+1}}{\uinst{\Renv}{\Penv \wedge \vv = \expr'}}
\end{align}
\end{proof}


\mypara{Completeness of Bounded Unfolding}
%
Next, we use the fact that unfolding preserves
equational links to show that bounded unfolding
is \emph{complete} for equational proofs.
%
That is, we prove by induction on the structure
of the equational proof that whenever there is
an \emph{equational proof} of $\expr = \expr'$,
there exists some bounded unfolding that
suffices to prove the equality.

\begin{lemma}[\textbf{Bounded Completeness}]\label{lem:bounded-completeness}
  If   $\eqpf{\Renv}{\Penv}{\expr}{\expr'}$
  then $\exists 0 \leq n.\ \smtIsValid{\binst{\Renv}{\Penv \wedge \vv = \expr}{n}}{\vv = \expr'}$.
\end{lemma}

\begin{proof} (Of Lemma~\ref{lem:bounded-completeness})
%
The proof follows by induction on the structure
of $\eqpf{\Renv}{\Penv}{\expr}{\expr'}$.
%
\smallskip \emph{Base Case \rulename{Eq-Refl}:}
Follows immediately as $\expr \equiv \expr'$.
%
\smallskip \emph{Inductive Case \rulename{Eq-Trans}}
\begin{align}
\intertext{In this case, there exists $\expr''$ such that}
  & \eqpf{\Renv}{\Penv}{\expr}{\expr''}
    \label{eq:bc:1} \\
  & \smtIsValid{\uinst{\Renv}{\Penv \wedge \vv = \expr''}}{\vv = \expr'}
    \label{eq:bc:2} \\
\intertext{By the induction hypothesis (\ref{eq:bc:1}) implies there exists $0 \leq n$ such that}
  & \smtIsValid{\binst{\Renv}{\Penv \wedge \vv = \expr}{n}}{\vv = \expr''}
    \notag \\
\intertext{By Lemma~\ref{lem:link} we have}
  & \smtIsValid{\binst{\Renv}{\Penv, \vv = \expr}{n+1}}{\uinst{\Renv}{\Penv \wedge \vv = \expr''}}
    \notag \\
\intertext{Thus, by (\ref{eq:bc:2}) and modus ponens we get}
  & \smtIsValid{\binst{\Renv}{\Penv, \vv = \expr}{n+1}}{\vv = \expr'} \notag
\end{align}
\end{proof}


\mypara{\pbesym is a Fixpoint of Bounded Unfolding}
%
Next, we show that the proof search procedure \pbesym computes
the least-fixpoint of the bounded unfolding and hence,
returns \ttrue iff there exists \emph{some} unfolding
depth $n$ at which the goal can be proved.

\begin{lemma}[\textbf{Fixpoint}]\label{lem:fixpoint}
$\pbe{\Renv}{\Penv}{\expr = \expr'}$ iff
$\exists n.\ \smtIsValid{\binst{\Renv}{\Penv, \vv = \expr}{n}}{\vv = \expr'}$
\end{lemma}

\begin{proof} (Of Lemma~\ref{lem:fixpoint})
%
Let $\Penv' \doteq \Penv \wedge \vv = \expr$.
%
The proof follows by observing that
$\pbe{\Renv}{\Penv}{\expr = \expr'}$
computes the \emph{least-fixpoint}
of the sequence
%
\begin{equation}
\Penv_i \doteq \binst{\Renv}{\Penv'}{i} \label{eq:fix:1}
\end{equation}
%
Specifically, we can prove by induction on $i$
that at each invocation of $\pbeloop{i}{\Penv_i}$
in Figure~\ref{fig:pbe}, $\Penv_i$ is equal to
$\binst{\Renv}{\Penv'}{i}$.

\emph{Case $\Rightarrow$:}
Assume that $\pbe{\Renv}{\Penv}{\expr = \expr'}$.
That is, at some iteration $i$ we have
$\smtIsValid{\Penv_i}{\vv = \expr'}$,
\ie by (\ref{eq:fix:1}) we have
$\smtIsValid{\binst{\Renv}{\Penv'}{i}}{\vv = \expr'}$.

\emph{Case $\Leftarrow$:}
Pick the smallest $n$ such that
$\smtIsValid{\binst{\Renv}{\Penv'}{n}}{\vv = \expr'}$.
Using Lemmas~\ref{lem:bounded:mono}~and~\ref{lem:bounded:fix}
we can then show that forall $0 \leq k < n$, we have
%
$$
\binst{\Renv}{\Penv'}{k} \subset \binst{\Renv}{\Penv'}{k+1}
$$
%
and
%
$$\binst{\Renv}{\Penv'}{k} \not \vdash {\vv = \expr'}$$
%
Hence, after $n$ iterations of the recursive loop,
$\pbe{\Renv}{\Penv}{\expr = \expr'}$, returns $\ttrue$.
\end{proof}

\mypara{Completeness of \pbesym}
%
Finally, by combining Lemma~\ref{lem:fixpoint}
and  Lemma~\ref{lem:link}  we can show that
\pbesym is complete, \ie if there is an equational proof
that $\expr = \expr'$ under $\Renv, \Penv$,
then \pbe{\Renv}{\Penv}{\expr}{\expr'} returns $\ttrue$.

\begin{theorem}[\textbf{Completeness}] \label{thm:completeness}
  \pbe{\Renv}{\Penv}{\expr = \expr'} = \ttrue iff
  \eqpf{\Renv}{\Penv}{\expr}{\expr'}
\end{theorem}

\subsection{\pbesym Terminates} \label{sec:pbe:termination}

So far, we have shown that our proof search
procedure $\pbesym$ is both sound and complete.
%
Of course, both of these are easy to achieve
simply by \emph{enumerating} all possible
instances and repeatedly querying the SMT
solver.
%
Of course, such an approach --- lets call it
``monkeys-with-typewriters'' --- is rather
impractical: it may never terminate.
%
Fortunately, next, we show that in addition
to being sound and complete, our proof search
procedure always terminates.
%
The key insight is that if $\pbesym$ \emph{does not terminate}
then it is because
%
(1)~the set of depth $n$ unfoldings form a strictly
    increasing infinite chain,
    $\binst{\Renv}{\Penv}{0} \subset \binst{\Renv}{\Penv}{1} \ldots$,
    which means that
(2)~there is an infinite sequence of application terms
    $\stepsmany{\fapp{\rfun_0}{\expr_0}}{\stepsmany{\fapp{\rfun_1}{\expr_1}}{\ldots}}$
    which means that
(3)~the reflected function $\rfun_0$ is \emph{not total}, contradicting the
    totality assumption (\ref{asm:total}).
%
Thus, \pbesym must terminate!
%
Next, we formalize the above intuition and
prove the termination Theorem~\ref{thm:termination}.

\mypara{Symbolic Step}
%
A pair
$\symstep{\fapp{\rfun}{\expr}}{\fapp{\rfun'}{\expr'}}$
is a
$\Renv,\Penv$-\emph{symbolic step} (abbrev. step)
if
%
\begin{itemize}
  \item $\rlookup{\Renv}{\rfun} \equiv \rdef{x}{\pred}{\body}$,
  \item $\smtIsValid{\Penv \wedge Q}{\pred_i}$ for some $\Renv$-instance $Q$,
  \item $\issubterm{\fapp{\rfun'}{\expr'}}{\SUBSTS{\body_i}{x}{\expr}}$
\end{itemize}

\mypara{Steps and Reductions}
%
Every symbolic step under an inhabited $\Penv$ corresponds
to a \emph{sequence} of steps in the concrete semantics:

\begin{lemma}[\textbf{Step-Reductions}]\label{lem:step-reductions}
If   $\symstep{\fapp{\rfun}{\expr}}{\fapp{\rfun'}{\expr'}}$
     is a symbolic step under $\Renv, \Penv$
     and
     $\satisfies{\Renv}{\store}{\Penv}$
then $\trans{\Renv}
            {\fapp{\rfun}{\apply{\store}{\expr}}}
            {\ctxapp{\ctx}{\fapp{\rfun'}{\apply{\store}{\expr'}}}}$
for some context $\ctx$.
\end{lemma}

\begin{proof} (Of Lemma~\ref{lem:step-reductions})
Using Lemmas~\ref{lem:subterm-eval}, \ref{lem:smt-approx}
and the definition of symbolic steps.
\end{proof}

\mypara{Steps and Values}
%
The following lemma states that if
$\symstep{\fapp{\rfun}{y}}{\expr'}$
is a symbolic step under a $\Penv$
inhabited by $\store$ then
$\apply{\store}{\fapp{\rfun}{y}}$
reduces to a value.
%
The proof follows by observing that if $\Penv$
is inhabited by $\store$, and a particular step
is possible, then the guard corresponding to that
step must also be true under $\store$ and hence,
by totality, the function must reduce to a value
under the given store.

\begin{lemma}[\textbf{Step-Value}]\label{lem:step-value}
If $\Renv$ is total,
   $\satisfies{\Renv}{\store}{\Penv}$, and
   $\symstep{\fapp{\rfun}{y}}{\expr'}$ is a $\Renv,\Penv$ step
then $\trans{\Renv}{\apply{\store}{\fapp{\rfun}{y}}}{\val}$.
\end{lemma}


\begin{proof} (Of Lemma~\ref{lem:step-value})
\begin{alignat}{3}
\bcoz{Assume that}
  && \Renv, \store & \models \Penv \label{eq:stepval:1} \\
\bcoz{Let}
  && \store^*               & \doteq \store \mapsingle{\params{x}}{\apply{\store}{\params{y}}} \label{eq:stepval:2}\\
  && \rlookup{\Renv}{\rfun} & \doteq \rdef{x}{\pred}{\body} \label{eq:stepval:3} \\
\intertext{As $\symstep{\fapp{\rfun}{y}}{\expr'}$ is a $\Renv,\Penv$ step,
for some $i$, \Renv-instance $Q$ we have}
  && \Penv \wedge Q & \vdash \SUBSTS{\pred_i}{x}{y} \notag \\
\bcoz{Hence, by (\ref{eq:stepval:1} and Lemma~\ref{lem:smt-approx}}
  && \Renv, \store  & \models \SUBSTS{\pred_i}{x}{y} \label{eq:stepval:4}\\
\bcoz{As}
  && \apply{\store^*}{\pred_i}
  & \equiv \apply{\store}{\SUBSTS{\pred_i}{x}{y}} \\
\bcoz{The fact (\ref{eq:stepval:4}) yields}
  && \Renv, \store^*  & \models \pred_i \notag \\
\bcoz{By the Totality Assumption~\ref{asm:total}}
  && \Renv & \models \stepsmany{\apply{\store^*}{\fapp{\rfun}{x}}}{\val} \notag \\
\bcoz{That is}
  && \Renv & \models \stepsmany{\apply{\store}{\fapp{\rfun}{y}}}{\val} \notag
\end{alignat}
\end{proof}

\mypara{Symbolic Trace}
%
A sequence
%
$\fapp{\rfun_0}{\expr_0},\fapp{\rfun_1}{\expr_1},\fapp{\rfun_2}{\expr_2},\ldots$
%
is a $\Renv,\Penv$-\emph{symbolic trace} (abbrev. trace)
if $\params{\expr_0} \equiv \params{x_0}$ for some variables $x_0$,
and $\symstep{\fapp{\rfun_i}{\expr_i}{\fapp{\rfun_{i+1}}{\expr_{i+1}}}}$
is a $\Renv,\Penv$-step, for each $i$.

\begin{lemma}[\textbf{Finite-Trace}]\label{lem:finite-trace}
If $\Renv$ is total, $\satisfies{\Renv}{\store}{\Penv}$
then every $\Renv,\Penv$-trace is finite.
\end{lemma}

\begin{proof} (Of Lemma~\ref{lem:finite-trace})
We prove this theorem by contradiction.
Specifically, we show that an infinite
$\Renv,\Penv$-trace for a $\Penv$
inhabited by $\store$ yields a
contradiction, and hence, all the
$\Renv, \Penv$ symbolic traces
must be finite.
%
\begin{alignat}{3}
\bcoz{Assume that}
   && \Renv, \store & \models \Penv
   \label{eq:fin:0}\\
\bcoz{and assume an infinite $\Renv, \Penv$ trace:}
   && \fapp{f_0}{\expr_0}, & \fapp{f_1}{\expr_1}, \ldots
   \label{eq:fin:1}\\
\bcoz{Where additionally}
   && \params{\expr_0} & \equiv \params{x_0}
      \label{eq:fin:2} \\
\bcoz{Define}
   && \expr^*_i        & \equiv \apply{\store}{\fapp{\rfun_i}{\expr_i}}
      \label{eq:fin:3} \\
\bcoz{By Lemma~\ref{lem:step-reductions}, for every $i \in \nats$}
   && \Renv & \models \stepsmany{\expr^*_i}{\ctxapp{\ctx_i}{\expr^*_{i+1}}}
      \notag \\
\bcoz{Hence, by Lemma~\ref{lem:div}}
   && \expr^*_0
    & \ \mbox{diverges under $\Renv$}
      \notag \\
\bcoz{\ie, by (\ref{eq:fin:2}, \ref{eq:fin:3})}
   && \apply{\store}{\fapp{\rfun_0}{x_0}}
    & \ \mbox{diverges under $\Renv$}
      \label{eq:fin:5} \\
\bcoz{But by (\ref{eq:fin:0}) and Lemma~\ref{lem:step-value}}
   && \Renv & \models \stepsmany{\apply{\store}{\fapp{\rfun_0}{x_0}}}
                                {\val}
      \quad \mbox{contradicting (\ref{eq:fin:5})} \notag
\end{alignat}
%
Hence, assumption (\ref{eq:fin:1}) cannot hold, \ie all
all the $\Renv, \Penv$ symbolic traces must be finite.
\end{proof}

\mypara{Unfolding Chains and Traces}
%
If the unfolding $\Renv, \Penv$ yields
an infinite chain, then $\Renv, \Penv$
has an infinite symbolic trace:

\begin{lemma} [\textbf{Ascending Chains}] \label{lem:chain}
Let $\Penv_i \doteq \binst{\Renv}{\Penv}{i}$.
If there exists an (infinite) ascending chain
  $$\Penv_0 \subset \ldots \subset \Penv_n \ldots$$
then there exists an (infinite) symbolic trace
  $$\fapp{\rfun_0}{\expr_0},\ldots,\fapp{\rfun_n}{\expr_n},\ldots$$
\end{lemma}

To prove the above,
we construct a trace
by selecting at level $i$
(\ie in $\Penv_i$),
an application term
$\fapp{\rfun_i}{\expr_i}$
that was created by
unfolding an application
term at level $i-1$
(\ie in $\Penv_{i-1}$).


\mypara{Transparency}
%
We say $\Penv$ is \emph{inconsistent}
if $\smtIsValid{\Penv}{\tfalse}$.
%
We say $\Penv$ is \emph{transparent}
if it is either inhabited or inconsistent.
%
As an example of a \emph{non-transparent}
$\Penv_0$ consider the predicate
$\kw{lenA xs} = 1 + \kw{lenB xs}$,
where \kw{lenA} and \kw{lenB} are both
identical definitions of the list length
function.
%
Clearly there is no $\store$ that causes
the above predicate to evaluate to $\ttrue$.
%
At the same time, the SMT solver cannot
(using the decidable, quantifier-free theories)
prove a contradiction as that requires
induction over \kw{xs}.


\mypara{\pbesym Terminates}
%
Finally, we can prove that our proof search
procedure $\pbesym$ terminates.

\begin{theorem}[\textbf{Termination}] \label{thm:termination}
  If $\Penv$ is transparent
  then $\pbe{\Renv}{\Penv}{\pred}$ terminates.
\end{theorem}

% \RJ{MAYBE KEEP THIS ONE}

\begin{proof} (Of Theorem~\ref{thm:termination})
As $\Penv$ is transparent, there are two cases.
%
\begin{alignat}{3}
%
\intertext{\quad \emph{Case: $\Penv$ is inconsistent.}}
%
\bcoz{By definition of inconsitency}
  && \Penv & \vdash \tfalse \notag \\
\bcoz{Hence}
  && \Penv & \vdash \pred \quad \mbox{\ie $\pbe{\Renv}{\Penv}{\pred}$ terminates immediately.}
     \notag \\%[0.1in]
%
\intertext{\quad \emph{Case: $\Penv$ is inhabited.}}
%
\bcoz{That is, exists $\store$ s.t.}
  && \Renv, \store & \models \Penv \label{eq:term:1} \\
\intertext{\quad Suppose that $\pbe{\Renv}{\Penv}{\pred}$ does not terminate.}
\bcoz{That is, there is an infinitely increasing chain:}
  && \Penv_0 \subset \ldots \subset \Penv_n \ldots
  \label{eq:term:2} \\
\bcoz{By Lemma~\ref{lem:chain}}
  && \Renv,\Penv & \ \mbox{has an infinite trace} \notag
\end{alignat}
Which, by (\ref{eq:term:1}) contradicts Lemma~\ref{lem:finite-trace}.
Thus, (\ref{eq:term:2}) is impossible, \ie $\pbe{\Renv}{\Penv}{\pred}$
terminates.
\end{proof}
